\chapter{Pendahuluan}

\section{Latar Belakang}
Stabilitas pasar keuangan modern seringkali terdistorsi oleh fenomena korelasi ekstrem antar aset yang meningkat drastis selama periode krisis ekonomi, sebuah kondisi yang secara teknis diidentifikasi sebagai \textit{Financial Dilemma}. Model \textit{Mean-Variance} (MV) yang diusulkan oleh \cite{abdul_hali_markowitz_2020} telah menjadi fondasi utama dalam \textit{Modern Portfolio Theory}, namun model ini memiliki keterbatasan krusial dalam mengasumsikan distribusi return yang normal dan stasioner \cite{abdul_hali_markowitz_2020}. Ketidakmampuan model MV untuk menangkap anomali pasar dan ketergantungan non-linear menyebabkan estimasi risiko menjadi tidak akurat, terutama ketika data historis tidak lagi mencerminkan dinamika masa depan secara reliabel \cite{sikalo_efficient_2022}. Oleh karena itu, diperlukan transformasi fungsional objektif portofolio ke dalam kerangka \textit{Quadratic Unconstrained Binary Optimization} (QUBO) agar kompleksitas tersebut dapat dipetakan secara matematis ke dalam arsitektur komputasi kuantum \cite{wang_achieving_2025}.

Dalam upaya menangkap interaksi strategis antar partisipan pasar, penggunaan \textit{Game Theory} menawarkan pendekatan yang lebih dinamis dibandingkan sekadar ekstrapolasi statistik imbal hasil. Integrasi kerangka kerja multi-agen ke dalam model \textit{Ising} bertujuan untuk menggantikan parameter medan lokal ($h_i$) konvensional dengan nilai \textit{marginal payoff} yang diturunkan dari utilitas strategis partisipan \cite{la_torre_multi-agent_2026}. Pendekatan teori permainan ini memungkinkan deteksi bias aset secara lebih presisi, di mana setiap aset dipandang sebagai pemain yang bereaksi terhadap kondisi pasar guna mencapai ekuilibrium \cite{sikalo_efficient_2022}. Dengan mengadopsi formulasi \textit{minimax} atau strategi kooperatif dalam penentuan bias, model Hamiltonian \textit{Ising} dapat merepresentasikan preferensi risiko investor secara lebih eksplisit dan adaptif terhadap guncangan pasar eksternal \cite{wei_solving_2025}.

Selain bias individu, pemodelan korelasi antar aset menuntut metrik yang mampu melampaui limitasi koefisien kovarians linear yang sering kali menyesatkan pada kondisi \textit{bear market}. \textit{Quantum Mutual Information} (QMI) diimplementasikan sebagai instrumen untuk menyusun koefisien kopling ($J_{ij}$) dalam sistem Hamiltonian guna menangkap seluruh spektrum korelasi probabilistik, baik yang bersifat linear maupun non-linear \cite{wang_achieving_2025}. Pemanfaatan QMI memungkinkan dekomposisi struktur interaksi aset menjadi energi interaksi mekanika kuantum yang lebih akurat, yang merupakan prasyarat mutlak bagi keberhasilan algoritma \textit{Variational Quantum Eigensolver} (VQE) \cite{wei_solving_2025}. Melalui sinergi antara \textit{Game Theory} dan QMI, optimasi portofolio tidak lagi hanya bergantung pada metrik statistik murni, melainkan pada sintesis informasi kuantum dan interaksi strategis yang lebih representatif terhadap realitas kompleksitas finansial.

Implementasi algoritma \textit{Variational Quantum Eigensolver} (VQE) merupakan terobosan signifikan yang mengadaptasi prinsip simulasi sistem kuantum untuk penyelesaian masalah optimasi diskrit di sektor finansial. Secara historis, VQE dikembangkan untuk mengestimasi \textit{ground state} pada sistem molekuler melalui minimisasi nilai ekspektasi Hamiltonian kimia yang bersifat \textit{non-commuting} \cite{fedorov_vqe_2022}. Namun, transformasi fungsional portofolio ke dalam model \textit{Ising} memungkinkan VQE digunakan untuk memetakan rezim pasar yang optimal melalui pencarian konfigurasi energi terendah dalam ruang pencarian QUBO yang bersifat \textit{NP-hard} \cite{wang_achieving_2025}. Fleksibilitas VQE dalam menangani perangkat \textit{Noisy Intermediate-Scale Quantum} (NISQ) menjadikannya instrumen yang lebih unggul dibandingkan algoritma kuantum murni seperti \textit{Quantum Phase Estimation} (QPE) yang membutuhkan kedalaman sirkuit yang tidak praktis untuk saat ini \cite{wang_achieving_2025}.

Metode optimasi iteratif konvensional dalam kerangka kerja VQE umumnya berlandaskan pada algoritma \textit{Gradient Descent} (GD) yang memperbarui parameter secara bertahap menuju arah negatif gradien fungsional biaya \cite{wang_achieving_2025}. Meskipun GD merupakan standar industri dalam pembelajaran mesin klasik, implementasinya pada perangkat keras kuantum menghadapi hambatan \textit{sampling overhead} yang sangat besar karena memerlukan evaluasi gradien untuk setiap dimensi parameter secara eksplisit melalui teknik seperti \textit{parameter-shift rule} \cite{dao_exploring_2025}. Sebagai solusi, algoritma \textit{Simultaneous Perturbation Stochastic Approximation} (SPSA) lebih dipilih dalam lingkungan \textit{Noisy Intermediate-Scale Quantum} (NISQ) karena kemampuannya melakukan efisiensi estimasi gradien hanya dengan dua pengukuran fungsional per iterasi, terlepas dari jumlah dimensi parameter yang ada \cite{wang_achieving_2025}. Selain efisiensi jumlah \textit{shot}, SPSA memiliki ketangguhan intrinsik terhadap fluktuasi statistik dan \textit{noise} pengukuran yang sering kali menyebabkan instabilitas pada algoritma berbasis gradien analitik \cite{dao_exploring_2025}.

Sinergi komputasi kuantum tersebut juga menuntut pemilihan \textit{ansatz} yang optimal, di mana kerangka kerja \textit{TwoLocal} sering kali menjadi fondasi dalam membangun sirkuit rotasi dan \textit{entanglement} yang repetitif \cite{dao_exploring_2025}. Meskipun \textit{TwoLocal} menawarkan fleksibilitas tinggi melalui kombinasi gerbang rotasi generik dan pola konektivitas yang luas, penggunaan \textit{ansatz} ini tanpa optimasi spesifik dapat menghasilkan jumlah gerbang dua-qubit yang berlebihan serta parameter redundan yang menyulitkan konvergensi \cite{wang_achieving_2025}. \textit{Ansatz EfficientSU(2)} hadir sebagai spesialisasi dari \textit{TwoLocal} yang secara eksplisit mengoptimalkan struktur gerbang $SU(2)$—seperti rotasi $RY$ dan $RZ$—serta gerbang \textit{CNOT} dalam pola yang lebih ringkas dan sesuai dengan topologi perangkat keras \cite{dao_exploring_2025}. Keunggulan utama \textit{EfficientSU(2)} dibandingkan \textit{TwoLocal} standar terletak pada kemampuannya mereduksi kedalaman sirkuit secara signifikan tanpa mengorbankan \textit{expressivity} yang dibutuhkan untuk memetakan dinamika portofolio finansial, sehingga meminimalkan efek dekoherensi selama proses optimasi \cite{dao_exploring_2025}.

\section{Rumusan Masalah}
Berdasarkan latar belakang tersebut, permasalahan yang akan dibahas dalam tugas akhir ini adalah:
\begin{enumerate}
    \item Bagaimana efek perumusan bias aset sebagai parameter medan lokal ($h_i$) dalam model Hamiltonian \textit{Ising} melalui pendekatan utilitas \textit{Game Theory} pada kinerja algoritma?
    \item Bagaimana mengonstruksi koefisien kopling ($J_{ij}$) antar aset menggunakan \textit{Quantum Mutual Information} (QMI) untuk menangkap interaksi non-linear?
    \item Bagaimana efektivitas algoritma \textit{Variational Quantum Eigensolver} (VQE) dengan optimasi spsa dan \textit{Ansatz EfficientSU(2)}  dalam mencari solusi keadaan dasar (\textit{ground state}) untuk menentukan keputusan alokasi aset yang optimal?
\end{enumerate}

\section{Tujuan Penelitian}
Dari rumusan masalah yang diberikan, tujuan dari pembuatan tugas akhir ini antara lain:
\begin{enumerate}
    \item Menganalisis formulasi bias aset sebagai parameter medan lokal ($h_i$) dalam model Hamiltonian \textit{Ising} berbasis kerangka kerja \textit{Game Theory} terhadap kinerja algoritma.
    \item Mengonstruksi koefisien kopling ($J_{ij}$) antar aset melalui pemanfaatan \textit{Quantum Mutual Information} (QMI) sebagai representasi interaksi non-linear.
    \item Mengevaluasi performa dan efektivitas algoritma \textit{Variational Quantum Eigensolver} (VQE) dengan optimasi spsa dan \textit{Ansatz EfficientSU(2)} dalam menemukan solusi optimal pada model portofolio kuantum yang telah dibangun.
\end{enumerate}

\section{Batasan Masalah}
Agar penelitian ini lebih terfokus, maka pengerjaan tugas akhir ini diberikan batasan-batasan sebagai berikut:
\begin{enumerate}
    \item Portofolio yang dianalisis dibatasi pada pemilihan 2 aset dari 4 pilihan aset yang berbeda.
    \item Model permainan yang digunakan adalah skema \textit{non zero-sum game}.
    \item Simulasi dan implementasi algoritma dilakukan menggunakan simulator \textit{Pennylane}.
\end{enumerate}

\section{Manfaat Penelitian}
Hasil dari tugas akhir ini diharapkan dapat memberikan alternatif metode optimasi portofolio yang lebih tangguh terhadap dinamika korelasi pasar serta memberikan kontribusi akademis dalam penerapan komputasi kuantum pada bidang finansial.

\newpage
\thispagestyle{empty}
\vspace*{\fill}
    \begin{center}
	Halaman ini sengaja dikosongkan
    \end{center}
\vspace*{\fill}
\pagebreak
